\chapter{Mesh Deformation}\label{ch:deform}

Deformation creates new shapes by editing existing ones: a ball that squashes when it bounces, a
character that walks or talks, a design variant obtained by pulling part of a model. Behind an
intuitive interface (the user grabs a handle and drags it) hide the mathematical concepts of this
chapter: interpolation and generalized barycentric coordinates, skinning, and variational methods
that minimise an energy measuring how much the surface is distorted. It follows Tema 11 and the
survey of \citet{Botsch2008}. We use the cotangent matrix $\mC$ of \cref{ch:smooth}.

\section{Controlling a deformation}

Moving vertices one by one is easy but tedious; moving a group of vertices rigidly is easy but too
restrictive, since it tears the surface at the border of the group. A practical interface lets
the user move a \emph{handle} region $H$ while keeping a region $F$ fixed, and the rest of the
surface, the \emph{region of interest}, must deform in an intuitive way. The simplest way of
propagating the motion is \emph{attenuation}: every vertex receives the handle displacement
$\bd$ scaled by a smooth falloff of its distance to the handle,
\begin{equation}
\bx'_i=\bx_i+\phi\bigl(\operatorname{dist}(\bx_i,H)/R\bigr)\,\bd,\qquad
\phi(s)=(1-s^2)^2\ \text{for}\ s<1,\ 0\ \text{otherwise}.
\end{equation}
The choice of distance is a trade-off between quality and computational cost: Euclidean distance
is cheap but leaks the deformation to parts that are close in space but far on the surface; the
minimal number of edges (a breadth-first search) is cheap and intrinsic but depends on the
tessellation; geodesic distance is intrinsic and accurate but more expensive (\cref{fig:attenuation}).
Attenuation translates vertices but never rotates them, so details are sheared rather than
rotated; the methods below solve this.

\fig[0.85\textwidth]{ch13_attenuation.png}{Moving a front hoof of the horse with attenuation
($R=0.3$, unit diagonal). Left: geodesic distance to the handle. Centre: with Euclidean distance the
displacement leaks to the chest and the other front leg. Right: with geodesic distance it stays on
the dragged leg.}{fig:attenuation}

Deformation methods can be classified as \emph{space deformations}, which deform the ambient
space and with it any embedded object (lattices, cages, skeletons), and \emph{surface-based}
methods, which optimise the surface directly (shell-based deformation on polygon meshes,
volumetric deformation on tetrahedral meshes). Multiscale methods split a shape into a smooth base
and details, deform the base, and add the details back, for efficiency and detail preservation.

\section{Interpolation}

Space deformations rest on interpolation inside cells. In one dimension, a value between $p_0$ and
$p_1$ at parameter $t\in[0,1]$ is $\operatorname{lerp}(p_0,p_1,t)=(1-t)p_0+tp_1$. \emph{Bilinear}
interpolation in a square cell interpolates along one direction and then along the other,
\begin{equation}
p(s,t)=(1-s)(1-t)\,p_{00}+s(1-t)\,p_{10}+(1-s)t\,p_{01}+st\,p_{11},
\end{equation}
and \emph{trilinear} interpolation in a cube does the same with eight corners. Inside a triangle, a
point is a unique convex combination of the corners, $\bx=\sum_i\lambda_i\bv_i$ with
$\lambda_i=A_i/A$, the ratios of the areas of the sub-triangles opposite to each corner
(\emph{barycentric coordinates}).

\paragraph{Generalized barycentric coordinates.} For a polygon with vertices $\bv_1,\dots,\bv_n$ we
want weights $w_i(\bx)$ with $\sum_iw_i=1$ and $\sum_iw_i\bv_i=\bx$ (linear precision), that are
smooth inside, and ideally non-negative. For convex polygons, \emph{Wachspress coordinates}
\citep{Wachspress1975} are
\begin{equation}
w_i(\bx)\propto\frac{A(\bv_{i-1},\bv_i,\bv_{i+1})}{A(\bx,\bv_{i-1},\bv_i)\,A(\bx,\bv_i,\bv_{i+1})},
\end{equation}
rational functions that are positive inside convex polygons. For arbitrary polygons (and, in 3D,
closed triangle meshes \citep{Ju2005}), \emph{mean value coordinates} \citep{Floater2003} are
\begin{equation}
w_i(\bx)\propto\frac{\tan(\alpha_{i-1}/2)+\tan(\alpha_i/2)}{\norm{\bv_i-\bx}},
\end{equation}
where $\alpha_i$ is the angle at $\bx$ in the triangle $(\bx,\bv_i,\bv_{i+1})$. They are defined and
smooth everywhere, but become negative in non-convex polygons (\cref{fig:coords}). \emph{Harmonic
coordinates} \citep{Joshi2007}, obtained by solving Laplace's equation inside the cage, are positive
everywhere but need a numerical solve.

\fig[0.9\textwidth]{ch13_coordinates.pdf}{The coordinate function of one vertex (red dot) for
Wachspress and mean value coordinates on a convex pentagon, and for mean value coordinates on a
non-convex polygon, where the function takes negative values (blue, minimum $-0.174$) behind the
reflex vertex.}{fig:coords}

\section{Free-form deformation}

Free-form deformation \citep{Sederberg1986} embeds the object in a lattice of control points and
deforms space: when the user moves the control points, every point of the object moves with the
cell that contains it. For each vertex, its local coordinates relative to the lattice are computed
once (for example, a vertex with global coordinates $(25.6,14.7)$ in a grid of unit cells starting at
$(20,12)$ has local coordinates $(5.6,2.7)$: cell $(5,2)$ and parameters $(0.6,0.7)$); after the
lattice moves, the vertex is relocated by interpolating the moved corners with the same
parameters. Trilinear interpolation gives a continuous but only piecewise smooth deformation; the
original method uses a tensor-product Bernstein polynomial over the whole lattice,
\begin{equation}
\bx'(s,t,u)=\sum_{i,j,k}B_i^l(s)\,B_j^m(t)\,B_k^n(u)\,\bp_{ijk},
\end{equation}
which is smooth everywhere (\cref{fig:ffd}). The same idea works with triangular cells and
barycentric coordinates, convex polygonal cells with Wachspress coordinates, and arbitrary
polygonal or polyhedral \emph{cages} with mean value or harmonic coordinates: the object is
expressed once in the coordinates of the cage, and editing the few cage vertices deforms it.

\fig[0.7\textwidth]{ch13_ffd.pdf}{Free-form deformation in 2D. A shape is embedded in a $4\times4$
lattice of control points (left); moving the control points deforms the space and the shape with
it (right).}{fig:ffd}

\section{Skinning}

\emph{Skinning} is a variant of free-form deformation in which the control structure is placed
inside the shape: a skeleton of bones connected by joints. Every vertex is attached to one or
several bones with weights $w_{ij}\ge0$, $\sum_jw_{ij}=1$, usually painted by artists or computed
automatically from distances or heat diffusion. When the skeleton is animated (by changing the joint
angles), each bone $j$ has a rigid transformation $\mT_j$ relative to the rest pose, and the
vertices follow the blend
\begin{equation}
\bx'_i=\sum_jw_{ij}\,\mT_j\bx_i ,
\end{equation}
known as \emph{linear blend skinning} (LBS) or skeletal subspace deformation (SSD)
\citep{MagnenatThalmann1988}. It is very fast (it runs on the GPU for every frame of a game), but
blending rotation matrices linearly does not produce a rotation: near joints the blended matrix
shrinks, the volume collapses when bending, and twisting by about $180^\circ$ pinches the surface into
the infamous \emph{candy-wrapper} artefact (\cref{fig:lbs}). Dual quaternion skinning
\citep{Kavan2008} blends rigid motions in a space where the result remains rigid and avoids it.

\fig[0.7\textwidth]{ch13_lbs.png}{Linear blend skinning of a cylinder with two bones. Left: the
weight of the upper bone. Centre: bending the joint by $100^\circ$ loses volume on the inner side.
Right: twisting by $170^\circ$ collapses the surface (candy wrapper).}{fig:lbs}

\section{Surface-based deformation}

The challenge of intuitive deformation is that it requires both a global change and the
preservation of local detail. Surface-based methods formulate this as the minimisation of an
energy over the deformed positions, subject to the handle and fixed constraints.

\subsection{Shell-based deformation}

Physically, a thin shell resists \emph{stretching}, a change of the first fundamental form, and
\emph{bending}, a change of the second fundamental form. Linearising these energies in terms of a
displacement field $\bd$ gives the simple quadratic energy
\begin{equation}
E(\bd)=\int_Sk_s\norm{\grad\bd}^2+k_b\norm{\Lap\bd}^2\dd A ,
\end{equation}
whose minimisers satisfy the Euler--Lagrange equation $-k_s\Lap\bd+k_b\Lap^2\bd=\bzero$ in the region
of interest, with $\bd$ prescribed on the handle and zero on the fixed region. Pure stretching
($k_b=0$) gives harmonic displacements (membranes, with sharp creases at the constraints); pure
bending ($k_s=0$) gives biharmonic displacements (thin plates, with smooth transitions). The
discretisation uses the cotangent Laplacian: $\Lap\approx\mM^{-1}\mC$ and
$\Lap^2\approx\mM^{-1}\mC\mM^{-1}\mC$, and one sparse linear system is solved per coordinate
\citep{Botsch2008}.

\subsection{Differential coordinates}

Absolute coordinates are a poor representation for shape-preserving edits: moving one vertex
changes nothing else, and the local shape around a vertex is encoded only implicitly. The
\emph{differential} (Laplacian, or delta) coordinates of a vertex encapsulate its local shape
\citep{Sorkine2006}:
\begin{equation}
\bdelta_i=\bx_i-\frac{1}{d_i}\sum_{j\in\mathcal N(i)}\bx_j
\qquad\text{or, with cotangent weights,}\qquad
\bdelta_i=\sum_{j\in\mathcal N(i)}C_{ij}(\bx_j-\bx_i)\approx-2H_iA_i\bn_i .
\end{equation}
The uniform version is the vector from the average of the neighbours to the vertex; the cotangent
version is a more geometric weighting, invariant to the triangulation and based on angles, that
approximates the mean curvature normal (\cref{ch:diffgeo}) and preserves the geometry better. In
matrix form $\bdelta=\mL\mX$ for all vertices at once.

\paragraph{Reconstruction.} The deltas store the shape but all information about location is lost:
$\mL$ is not invertible, since adding a constant to all positions does not change $\mL\mX$ (its
rank is $n-1$ for a connected mesh). To recover the absolute coordinates the mesh must be placed
somewhere: specify the position of one vertex by appending a row with a one in its column (and
zeros elsewhere) and its position on the right-hand side. The extended system has a unique
solution, which reproduces the original mesh.

\subsection{Laplacian editing}

Laplacian editing \citep{Sorkine2004} combines two terms:
\begin{itemize}
  \item a \emph{match} term that moves the handle vertices as close as possible to the positions
  $\bc_k$ specified by the user (and keeps the fixed vertices in place);
  \item a \emph{regularisation} term that preserves the original local shape (angles, areas,
  lengths) through the differential coordinates.
\end{itemize}
Without regularisation only the handle moves; without the match term the mesh does not move at all;
with both, the result is a compromise:
\begin{equation}
\min_{\mX'}\ \norm{\mL\mX'-\bdelta}^2+w^2\sum_{k\in H\cup F}\norm{\bx'_k-\bc_k}^2 .
\end{equation}
Stacking the constraints below $\mL$ as rows of a selection matrix $\mS$ (scaled by $w$), the solution
satisfies the normal equations
\begin{equation}
\bigl(\mL\T\mL+w^2\mS\T\mS\bigr)\mX'=\mL\T\bdelta+w^2\mS\T\mathbf c ,
\end{equation}
a sparse symmetric positive definite system factorised once and solved for the three coordinates
and for every new handle position. The weight $w$ controls how strictly the constraints are
enforced (\cref{fig:soft}). With large $w$ they act as hard constraints, which can also be imposed
exactly by eliminating the constrained variables.

\fig[0.95\textwidth]{ch13_soft_weights.png}{Laplacian editing with soft constraints: the tail tip
of the cat is lifted by $0.2$ while the half of the body farthest from it is fixed. With $w=0.003$ the
regularisation dominates and the handle stays $0.083$ away from its target; with $w=0.03$ the error is
$0.002$; with $w=10$ the constraint is met exactly ($8\times10^{-6}$).}{fig:soft}

\paragraph{The rotation problem.} The differential coordinates are vectors expressed in the global
frame: they are invariant to translation, but not to rotation. When the handle rotates a part of the
surface, the reconstructed details keep their original orientation instead of rotating with the
surface, so they are sheared and distorted (\cref{fig:arap}, centre). \citet{Sorkine2004} estimate a
linearised rotation (plus scaling) per vertex as part of the same least-squares problem; the
non-linear method below solves the problem properly.

\subsection{As-rigid-as-possible deformation}

\citet{Sorkine2007} ask each one-ring to deform as rigidly as possible: there should be a rotation
$\mR_i$ per vertex that maps the original edges to the deformed ones,
\begin{equation}
E(\mX',\{\mR_i\})=\sum_i\sum_{j\in\mathcal N(i)}w_{ij}\,\bigl\|(\bx'_i-\bx'_j)-\mR_i(\bx_i-\bx_j)\bigr\|^2 ,
\end{equation}
with cotangent weights $w_{ij}$. The energy is minimised by alternating two steps:
\begin{enumerate}
  \item \textbf{Local step}: for fixed positions, each $\mR_i$ is the best rotation between the
  original and deformed edges of the one-ring, the Procrustes problem of \cref{ch:icp}: with
  $\mS_i=\sum_jw_{ij}\be_{ij}\be_{ij}'^{\mathsf T}=\mU_i\mSigma_i\mV_i\T$,
  $\mR_i=\mV_i\mU_i\T$ (with the determinant correction).
  \item \textbf{Global step}: for fixed rotations, the positions solve the Poisson-like system
  \begin{equation}
  \sum_{j\in\mathcal N(i)}w_{ij}(\bx'_i-\bx'_j)=\sum_{j\in\mathcal N(i)}\frac{w_{ij}}{2}(\mR_i+\mR_j)(\bx_i-\bx_j),
  \end{equation}
  that is $-\mC\mX'=\mathbf B$ with the constrained vertices moved to the right-hand side. The matrix
  does not change between iterations, so it is factorised once.
\end{enumerate}
Each step decreases the energy, and a few tens of iterations, initialised with the Laplacian editing
result, suffice in practice. Details now rotate with the surface (\cref{fig:arap}).

\fig[0.85\textwidth]{ch13_laplacian_arap.png}{A plane with bumps whose left edge is fixed (blue) and
whose right edge (red) is lifted and moved inwards. Laplacian editing keeps the bumps oriented
vertically, so they are sheared; as-rigid-as-possible deformation rotates them with the surface and
preserves their shape.}{fig:arap}

\begin{keypoints}
\begin{itemize}[leftmargin=*]
  \item Attenuation propagates handle displacements with a falloff; geodesic distances avoid leaks
  between nearby parts.
  \item Space deformations interpolate a lattice or cage (trilinear, Bernstein, Wachspress, mean
  value, harmonic coordinates); skinning blends bone transformations, with volume loss and candy
  wrappers.
  \item Shell-based deformation minimises linearised stretching and bending energies, leading to
  harmonic or biharmonic systems.
  \item Laplacian editing preserves differential coordinates $\mL\mX$ in the least-squares sense with
  soft handle constraints; ARAP adds per-vertex rotations, solved by local SVDs and a global Poisson
  system.
\end{itemize}
\end{keypoints}

\section*{Exercises}
\begin{exercise}
Show that mean value coordinates reproduce linear functions:
$\sum_iw_i(\bx)\bv_i=\bx$.
\end{exercise}
\begin{exercise}
Linear blend skinning of a point on the axis of a twisting cylinder: compute the blended position
for two bones rotated by $0$ and $\theta$ about the axis with weights $\tfrac12$ each, and show that the
distance of an off-axis point to the axis shrinks by $\abs{\cos(\theta/2)}$.
\end{exercise}
\begin{exercise}
Prove that the Laplacian matrix of a connected mesh has rank $n-1$ and that fixing one vertex makes
the reconstruction unique.
\end{exercise}
